\documentclass[12pt,a4paper]{article}
\usepackage[utf8]{inputenc}
\usepackage{amsmath,amsthm,amsfonts,amssymb}
\usepackage{ctex}
\usepackage{geometry}
\usepackage{fancyhdr}


% 页边距设置
\geometry{left=2cm, right=2cm, top=2.5cm, bottom=2.5cm}

% 页眉页脚设置
\pagestyle{fancy}
\fancyhf{}
\fancyhead[L]{数学试卷}
\fancyhead[R]{第 \thepage 页}

% 定理环境设置
\newtheorem{definition}{定义}
\newtheorem{theorem}{定理}
\newtheorem{lemma}{引理}
\newtheorem{corollary}{推论}

\begin{document}
	
	\title{2024数学强基数学分析II期末试题}
	\author{柯力洲}
	\date{\today}
	
	\maketitle
	
	\section*{定义}
	\begin{definition}
		$lebesgue$零测集
	\end{definition}
	
	\begin{definition}
	积分第二中值定理
	\end{definition}
	
		\begin{definition}
	$f_n(x)$一致收敛于$f(x)$
	\end{definition}

		\begin{definition}
	连通集
	\end{definition}
	
	\begin{definition}
		含参变量积分不一致收敛的$Cauchy$法则
	\end{definition}
	
	
	
	\section*{计算}
	\begin{enumerate}
		\item 计算下列表达式的值：
		\begin{equation*}
		\lim\limits_{x \to 0} \dfrac{\int_{0}^{x} \frac{1-\cot t}{t} dt}{x^2}
		\end{equation*}
		
		\item 求下列级数的收敛半径:
		\begin{equation*}
			\sum\limits_{n=1}^{\infty} n^nx^{n^2}
		\end{equation*}
		
		\item $\frac{x^2}{a^2}+\frac{y^2}{b^2}=1$在$x_0(\frac{a}{\sqrt{2}},\frac{b}{\sqrt{2}})$处的内法线方向为$U$，求$z=1-\frac{x^2}{a^2}-\frac{y^2}{b^2}$在$x_0$处$U$方向的导数
		
		\item 已知$x+y=y+v,xe^u=vy$求：
		\[
		\begin{bmatrix}
			\frac{\partial u}{\partial x} & \frac{\partial u}{\partial y} \\
			\frac{\partial v}{\partial x} & \frac{\partial v}{\partial y}
		\end{bmatrix}
		\]
		
		\item 求曲线$x^2+2y^2+3z^2=21$平行于$x+4y+6z=0$的切平面方程
		
		\item 对下列函数麦克劳林展开四项：
		\begin{equation*}
			f(x)=log(1+e^x)
		\end{equation*}
		
		\item 求$f_{xy}$,其中：
		\begin{equation*}
			f(x,y)=\int_{\frac{x}{y}}^{xy}(x-ty)\phi (t)dt
		\end{equation*}
		
		
	\end{enumerate}
	
	
	
	
	\section*{证明}
	\begin{enumerate}
		\item 证明下列定理：\\

			设$E \subset R^n \quad f:E \to R$对$\forall \delta >0$记$\Omega_f (\delta)=\sup\limits_{\vec{x},\vec{y} \in E ; |\vec{x}-\vec{y}|<\delta}|f(\vec{x})-f(\vec{y})|$
			则$f$在$E$上一致连续的充要条件是$\lim\limits_{\delta \to 0^+} \Omega_f (\delta)=0$

		
		\item $u_n(x)=\dfrac{(-1)^n(x+n)^n}{n^{n+1}}$,求$\sum\limits_{n+1}^{\infty} u_n(x)$的一致收敛性和绝对收敛性
		
		\item 求$\alpha$的范围使得$f(x,y)$可微
		\[
		f(x,y) =
		\begin{cases}
			(x^2+y^2)^{\alpha} \sin \dfrac{1}{x^2+y^2}, & \text{if } x^2+y^2 \neq 0 \\
			0, & \text{if } x^2+y^2 = 0 \\
		\end{cases}
		\]
		
		\item 求$I_n(\alpha)$其中：\\
		$I_n(\alpha)=\int_{0}^{+\infty} \dfrac{dx}{(x^2+a)^{n+1}}$
	\end{enumerate}
	
	
	
\end{document}
